\chapter{Electric Potential}

\section{Work}

\textbf{Work} (\(W\)) is the change in energy of a system
and is measured in joules (\unit{\joule}).

\section{Electric Potential Energy}

\textbf{Electric potential energy} (\(U_E\))
has SI units of joules (\unit{\joule}).
Electric potential energy can be found by integrating electric force
over distance,
i.e.,
\begin{equation}
    U_E = - \int \vecb{F_E} \cdot \odif{\vecb{r}}
\end{equation}

\section{Electric Potential}

\textbf{Electric potential} (\(V_E\) or simply \(V\))
has SI units of joules per coulomb (\unit{\joule \per \coulomb})
or volts (\unit{\volt}).
The difference in electric potential between two points
is also known as \textbf{voltage}.
Electric potential can be found be integrating electric field over distance,
i.e.,
\begin{equation}
    V = - \int \vecb{E} \cdot \odif{\vecb{r}}
\end{equation}
If we apply this to Coulomb's law [\pgRef{eq:coulombs-law}],
we find
\begin{gather}
    \begin{aligned}
        V &= \int \frac{k q}{r^2} \odif{r} \\
        &= \int - \frac{k q}{r} \odif{r}.
    \end{aligned}
\end{gather}

Similarly to how we can multiply electric field by charge
to find electric force,
we can multiply electric potential by charge to find electric potential energy.
